\section{Systems 扩展：GPU、并行和推理}

\subsection{Kernels}

GPU 的关键属性不是延迟，而是吞吐。kernel 设计的第一原则就是减少数据搬运，把高频数据留在更快的存储层级。

\begin{figure}[H]
\centering
\begin{tikzpicture}[every node/.style={font=\small, align=center}]
\node[draw, rounded corners, fill=gray!10, minimum width=3.0cm, minimum height=0.8cm] (dram) {DRAM\\warehouse};
\node[draw, rounded corners, fill=gray!20, right=2.0cm of dram, minimum width=3.0cm, minimum height=0.8cm] (sram) {SRAM\\factory floor};
\node[draw, rounded corners, fill=gray!30, below=1.2cm of sram, minimum width=3.4cm, minimum height=0.8cm] (comp) {compute units};
\draw[-{Latex[length=3mm]}, thick] (dram) -- node[above]{movement cost} (sram);
\draw[-{Latex[length=3mm]}, thick] (sram) -- (comp);
\end{tikzpicture}
\caption{GPU/内存的类比：把计算组织到更靠近数据的位置}
\end{figure}

\begin{warningbox}{kernel 设计的核心不是语法}
真正重要的是 memory access pattern、并行粒度和数据复用率。
\end{warningbox}

\subsection{Parallelism}

多 GPU 并行本质上是在更高一层重新做数据流设计。你可以切 batch、切 tensor、切 layer，也可以切 sequence。核心原则没变：少搬运、重同步、合理切分。

\begin{figure}[H]
\centering
\begin{tikzpicture}[node distance=1.0cm, every node/.style={font=\small, align=center}]
\node[draw, rounded corners, fill=blue!10, minimum width=2.4cm, minimum height=0.8cm] (g1) {GPU 1};
\node[draw, rounded corners, fill=blue!10, right=1.3cm of g1, minimum width=2.4cm, minimum height=0.8cm] (g2) {GPU 2};
\node[draw, rounded corners, fill=blue!10, below=1.1cm of g1, minimum width=2.4cm, minimum height=0.8cm] (g3) {GPU 3};
\node[draw, rounded corners, fill=blue!10, right=1.3cm of g3, minimum width=2.4cm, minimum height=0.8cm] (g4) {GPU 4};
\node[draw, rounded corners, fill=orange!18, below=1.4cm of g3, minimum width=5.4cm, minimum height=0.8cm] (ops) {collectives: gather / reduce / all-reduce};
\draw[-{Latex[length=3mm]}, thick] (g1) -- (g2);
\draw[-{Latex[length=3mm]}, thick] (g1) -- (g3);
\draw[-{Latex[length=3mm]}, thick] (g2) -- (g4);
\draw[-{Latex[length=3mm]}, thick] (g3) -- (g4);
\draw[-{Latex[length=3mm]}, thick] (g3) -- (ops);
\end{tikzpicture}
\caption{多 GPU 并行：参数、梯度、激活和 optimizer state 都能切开}
\end{figure}

\begin{table}[H]
\centering
\begin{tabular}{p{3cm}p{4.2cm}p{5.8cm}}
\toprule
\textbf{并行方式} & \textbf{切分对象} & \textbf{目的} \\
\midrule
Data parallelism & batch & 提升吞吐 \\
Tensor parallelism & matrix / attention tensors & 单层太大时拆开算 \\
Pipeline parallelism & layers & 把深层网络切段 \\
Sequence parallelism & sequence dimension & 支撑更长上下文 \\
\bottomrule
\end{tabular}
\caption{课程会反复出现的并行类型}
\end{table}

\subsection{Inference}

推理是实际用户反复付费的地方。训练是一次性的，推理却会在产品、RL 和评测中不断发生，因此全局上推理计算常常会超过训练。

\begin{figure}[H]
\centering
\begin{tikzpicture}[node distance=1.1cm, every node/.style={font=\small, align=center}]
\node[draw, rounded corners, fill=green!12, minimum width=3.2cm, minimum height=0.8cm] (pf) {prefill\\tokens known};
\node[draw, rounded corners, fill=green!22, right=1.8cm of pf, minimum width=3.2cm, minimum height=0.8cm] (dc) {decode\\one token at a time};
\node[draw, rounded corners, fill=green!32, below=1.1cm of dc, minimum width=4.5cm, minimum height=0.8cm] (opt) {KV cache / batching / speculative decoding};
\draw[-{Latex[length=3mm]}, thick] (pf) -- (dc);
\draw[-{Latex[length=3mm]}, thick] (dc) -- (opt);
\end{tikzpicture}
\caption{推理的两阶段：prefill 更像训练，decode 更像逐 token 生成}
\end{figure}

\begin{importantbox}{推理优化的目标}
让每一次生成都更便宜、更快、更稳定。
\end{importantbox}

